\documentclass{ctexart}
\usepackage{tikz}
\usepackage{expl3}

\begin{document}

\ExplSyntaxOn

\int_new:N \l_tmpc_int
\tl_new:N \l_tmpd_tl

% 定义颜色映射
\prop_new:N \g_color_map_prop
\prop_gput:Nnn \g_color_map_prop {g} {gray}
\prop_gput:Nnn \g_color_map_prop {b} {blue}
\prop_gput:Nnn \g_color_map_prop {r} {red}
\prop_gput:Nnn \g_color_map_prop {w} {white}

% 主函数：根据字符串绘制彩色正方形（九宫格排列）
\cs_new:Npn \draw_squares:n #1 {
    \begin{tikzpicture}[x=1cm, y=1cm]
        \int_zero:N \l_tmpa_int  % 用于计数
        
        % 按行优先顺序排列（第一行，然后第二行，最后第三行）
        \str_map_inline:nn {#1} {
            \int_incr:N \l_tmpa_int
            
            % 计算行和列（从0开始）
            \int_set:Nn \l_tmpb_int {\int_div_truncate:nn {\l_tmpa_int - 1} {3}}  % 行号 (0,1,2)
            \int_set:Nn \l_tmpc_int {\int_mod:nn {\l_tmpa_int - 1} {3}}           % 列号 (0,1,2)
            
            % 获取颜色
            \prop_get:NnN \g_color_map_prop {##1} \l_tmpd_tl
            
            % 绘制正方形（1cm边长，填充对应颜色，白色描边，2pt线宽）
            % 注意：Y坐标需要反转（第一行在顶部）
            \fill[\l_tmpd_tl, draw=white, line~width=2pt] 
                (\l_tmpc_int, 2 - \l_tmpb_int) rectangle ++(1, 1);
        }
    \end{tikzpicture}
}

% 创建一个更简洁的版本，不计算坐标，直接使用硬编码位置
% \cs_new:Npn \draw_squares_3x3_simple:n #1 {
%     \begin{tikzpicture}[x=1cm, y=1cm]
%         % 第一行 (y=2)
%         \prop_get:NnN \g_color_map_prop {\str_item:Nn {#1} {1}} \l_tmpa_tl
%         \fill[\l_tmpa_tl, draw=white, line~width=2pt] (0, 2) rectangle (1, 3);
        
%         \prop_get:NnN \g_color_map_prop {\str_item:Nn {#1} {2}} \l_tmpa_tl
%         \fill[\l_tmpa_tl, draw=white, line~width=2pt] (1, 2) rectangle (2, 3);
        
%         \prop_get:NnN \g_color_map_prop {\str_item:Nn {#1} {3}} \l_tmpa_tl
%         \fill[\l_tmpa_tl, draw=white, line~width=2pt] (2, 2) rectangle (3, 3);
        
%         % 第二行 (y=1)
%         \prop_get:NnN \g_color_map_prop {\str_item:Nn {#1} {4}} \l_tmpa_tl
%         \fill[\l_tmpa_tl, draw=white, line~width=2pt] (0, 1) rectangle (1, 2);
        
%         \prop_get:NnN \g_color_map_prop {\str_item:Nn {#1} {5}} \l_tmpa_tl
%         \fill[\l_tmpa_tl, draw=white, line~width=2pt] (1, 1) rectangle (2, 2);
        
%         \prop_get:NnN \g_color_map_prop {\str_item:Nn {#1} {6}} \l_tmpa_tl
%         \fill[\l_tmpa_tl, draw=white, line~width=2pt] (2, 1) rectangle (3, 2);
        
%         % 第三行 (y=0)
%         \prop_get:NnN \g_color_map_prop {\str_item:Nn {#1} {7}} \l_tmpa_tl
%         \fill[\l_tmpa_tl, draw=white, line~width=2pt] (0, 0) rectangle (1, 1);
        
%         \prop_get:NnN \g_color_map_prop {\str_item:Nn {#1} {8}} \l_tmpa_tl
%         \fill[\l_tmpa_tl, draw=white, line~width=2pt] (1, 0) rectangle (2, 1);
        
%         \prop_get:NnN \g_color_map_prop {\str_item:Nn {#1} {9}} \l_tmpa_tl
%         \fill[\l_tmpa_tl, draw=white, line~width=2pt] (2, 0) rectangle (3, 1);
%     \end{tikzpicture}
% }

\ExplSyntaxOff

% 使用示例
\section*{九宫格排列示例}

\subsection*{字符串 "gggrbwgrg" 对应的彩色正方形九宫格}

\noindent
\ExplSyntaxOn
\draw_squares:n {gggrbwgrg}
\ExplSyntaxOff

% 添加解释说明
\vspace{1cm}
\noindent
\textbf{字符串与九宫格对应关系：}
\begin{itemize}
    \item 字符串：\texttt{g g g r b w g r g}
    \item 九宫格位置：
    \begin{center}
    \begin{tabular}{|c|c|c|}
        \hline
        第1个(g) & 第2个(g) & 第3个(g) \\
        \hline
        第4个(r) & 第5个(b) & 第6个(w) \\
        \hline
        第7个(g) & 第8个(r) & 第9个(g) \\
        \hline
    \end{tabular}
    \end{center}
    \item 颜色映射：g=灰色，b=蓝色，r=红色，w=白色
\end{itemize}

\end{document}